% chapters/14/section-02.tex
\section{Animator dan Peralihan Animasi}
\label{sec:animator}

\subsection{Mesin Keadaan untuk Animasi}

Komponen \istilah{Animator} adalah mesin keadaan yang khusus mengurus animasi.
Gagasannya sama dengan mesin keadaan pada Bab~IX, tetapi yang berpindah bukan
perilaku benda, melainkan gambar yang ditampilkan.

Keadaan pada \istilah{Animator} disebut \istilah{animation clip}, yaitu
rangkaian gambar yang diputar berulang. Transisinya diatur melalui
\istilah{parameter}, yaitu nilai yang ditetapkan oleh kode Anda. Susunan untuk
pemain platformer diperlihatkan pada panel kiri
Gambar~\ref{fig:animator-canvas}.

\begin{figure}[htbp]
\centering
% chapters/14/figures/animator-dan-canvas.tikz
% Susunan dua hal yang dibahas pada bab ini: mesin keadaan animasi, dan
% antarmuka pengguna beserta susunan bendanya.
\begin{tikzpicture}[
  keadaan/.style={
    rectangle, rounded corners=8pt, draw=ubPrimary, thick, fill=ubAlt,
    align=center, text width=18mm, inner sep=4pt, minimum height=10mm,
    font=\footnotesize
  },
  panah/.style={-{Stealth[length=2.2mm]}, thick, draw=ubAccent},
  label/.style={font=\scriptsize, align=center},
  judul/.style={font=\footnotesize\bfseries, ubPrimary, align=center}
]

% ---------- Panel A: mesin keadaan animasi ----------
\begin{scope}
  \node[judul, anchor=south] at (2.5,1.5) {mesin keadaan animasi};

  \node[keadaan] (diam) at (0,0) {diam\\(idle)};
  \node[keadaan] (lari) at (5.0,0) {berlari\\(run)};
  \node[keadaan] (lompat) at (2.5,-2.0) {melompat\\(jump)};

  \draw[panah] (diam) -- node[label, above, yshift=2pt]
    {kecepatan $>$ 0,1} (lari);
  \draw[panah] (lari) to[bend left=20] node[label, above, yshift=1pt]
    {kecepatan $\approx$ 0} (diam);
  \draw[panah] (lari) -- node[label, right, xshift=3pt]
    {tanah = salah} (lompat);
  \draw[panah] (lompat) -- node[label, left, xshift=-3pt]
    {tanah = benar} (diam);
\end{scope}

% ---------- Panel B: susunan antarmuka ----------
\begin{scope}[xshift=9.4cm]
  \node[judul, anchor=south] at (2.0,1.5) {susunan antarmuka pengguna};

  \node[keadaan, text width=30mm] (canvas) at (2.0,0.6) {Canvas};
  \node[keadaan, text width=16mm] (skor) at (0.3,-1.0) {Teks Skor};
  \node[keadaan, text width=16mm] (nyawa) at (2.0,-1.0) {Teks Nyawa};
  \node[keadaan, text width=16mm] (ikon) at (3.7,-1.0) {Ikon Nyawa};

  \draw[panah] (canvas) -- (skor);
  \draw[panah] (canvas) -- (nyawa);
  \draw[panah] (canvas) -- (ikon);

  \node[label, anchor=north, text width=50mm] at (2.0,-1.7)
    {Seluruh elemen antarmuka ditempatkan di bawah satu \istilah{Canvas},
     sehingga susunannya tetap sama pada setiap ukuran layar.};
\end{scope}
\end{tikzpicture}

\caption{Mesin keadaan animasi dan susunan antarmuka pengguna}
\label{fig:animator-canvas}
\end{figure}

Jendela \istilah{Animator} dibuka melalui menu \istilah{Window}, lalu
\istilah{Animation}, lalu \istilah{Animator}. Di sana keadaan dan transisinya
digambar dengan cara yang sama seperti diagram keadaan pada
Gambar~\ref{fig:diagram-keadaan}: kotak untuk keadaan, dan anak panah untuk
transisi \cite{unity-leveldesign}.

\subsection{Jenis Parameter}

Tiga jenis parameter yang tersedia perlu Anda bedakan.

\begin{itemize}
  \item \textbf{Angka pecahan} (\istilah{float}), dipakai untuk nilai yang
        berubah halus, misalnya besar kecepatan.
  \item \textbf{Nilai benar atau salah} (\istilah{bool}), dipakai untuk keadaan
        yang hanya memiliki dua kemungkinan, misalnya sedang menyentuh tanah
        atau tidak.
  \item \textbf{Pemicu} (\istilah{trigger}), dipakai untuk kejadian yang terjadi
        sekali dan tidak memiliki lama, misalnya lompatan atau menerima
        kerusakan.
\end{itemize}

Pilihan yang tepat di antara ketiganya menghemat pekerjaan. Sebuah lompatan
sebaiknya memakai \istilah{trigger}, bukan \istilah{bool} yang harus
dikembalikan ke nilai semula oleh kode, sebab \istilah{trigger} dikembalikan
sendiri oleh mesin setelah dipakai.

\subsection{Menghubungkan Animasi dengan Keadaan Pemain}

Kode Anda tidak memutar animasinya sendiri, melainkan hanya menetapkan nilai
parameternya. Seluruh pekerjaan selebihnya dikerjakan mesin.

\begin{kodekeluaran}
animator.SetFloat(SpeedParam, Mathf.Abs(body.velocity.x));
animator.SetBool(GroundedParam, kontrol.IsOnGround);
\end{kodekeluaran}

\begin{peringatanbox}
Nama parameter yang dipakai pada kode harus sama persis dengan nama yang
dibuat pada jendela \istilah{Animator}, termasuk besar-kecil hurufnya.
Kesamaan itu \textbf{tidak diperiksa} oleh penyunting. Apabila namanya berbeda,
program tetap berjalan tanpa kesalahan, tetapi animasinya tidak pernah
berganti. Kebiasaan yang mencegahnya: pakai nama yang tetap, misalnya dengan
menyimpannya pada peubah tetap seperti pada skrip Daftar~\ref{lst:ch14-animasi}.
\end{peringatanbox}

\begin{tipsbox}
Animasi tidak selalu diperlukan. Pada permainan yang sederhana, gerak yang
halus sudah cukup memberi kesan hidup, sedangkan animasi menambah pekerjaan,
ukuran berkas, dan satu hal lagi yang dapat salah. Gunakanlah animasi pada hal
yang memang paling terlihat pemain, misalnya gerak pemain dan musuh.
\end{tipsbox}
